\documentclass[11pt]{article}
\usepackage[margin=1in]{geometry}
\usepackage{amsmath,amssymb,amsthm,booktabs,hyperref}
\newtheorem{theorem}{Theorem}
\newtheorem{lemma}{Lemma}
\newtheorem{proposition}{Proposition}
\title{Two refinements of the October 2026 OpenAI math release}
\author{AI-generated research draft prepared with Codex}
\date{October 9, 2026}
\begin{document}
\maketitle
\begin{abstract}
We derive a conditional quantitative refinement of the released van der
Waerden construction, and an explicit obstruction to improving the matrix
multiplication exponent using only its stated profile inequalities.
The former admits the coefficient $c=1/200$, versus the manuscript's
$10^{-5}$, subject to the upstream geometric and counting arguments.
The latter is a limitation of the abstract inequalities, not a lower bound
on the actual matrix multiplication exponent. These are research drafts:
neither literature priority nor independent formal verification is claimed.
\end{abstract}

\section{Scope and sources}
The source is \href{https://github.com/openai/math}{openai/math}, pinned at
\texttt{fd4aeeb2ee4fc729c18d98444fed42fd0529eeeb}.
The relevant papers are \emph{Quantitative Superexponential Bounds for van
der Waerden Numbers} (September 23, 2026) and \emph{An Upper Bound of
$9/4$ for the Matrix Multiplication Exponent} (October 2, 2026).
The inspected source revision and file SHA-256 hashes are recorded in
\texttt{source-manifest.json}. We use their constructions and identified
lemmas, not their headline theorems as black boxes for stronger conclusions.
Direct original sources are the
\href{https://github.com/openai/math/tree/fd4aeeb2ee4fc729c18d98444fed42fd0529eeeb/preprints/Quantitative-Superexponential-Bounds-for-van-der-Waerden-Numbers-September-23-2026}{van der Waerden manuscript}
and the
\href{https://github.com/openai/math/tree/fd4aeeb2ee4fc729c18d98444fed42fd0529eeeb/preprints/Matrix-Multiplication-Nine-Fourths-October-2-2026}{matrix multiplication manuscript}.
This independent follow-up is not an official OpenAI publication.
The parameter transfer below has been audited in source; it has not been
formalized in Lean. Finite algebra checks are supplementary evidence.

\section{A sharper norm-band multiplicity}
The van der Waerden construction uses keys comprising a coordinate box
and an integer squared-norm band. Its affinity and separation lemmas put
the points of a fixed key on a real affine line, at mutual distance at
least one. The released multiplicity lemma gives four occurrences.

\begin{lemma}\label{lem:band}
Let points on a real affine line be pairwise separated by distance at
least one. At most two of them can have squared norms in a single
half-open interval $[m,m+1)$.
\end{lemma}
\begin{proof}
Use signed arclength $t$ from the nearest point of the line to the origin.
The squared norm is $t^2+C$, for fixed $C\ge0$. If three points were
in the band, two of their parameters would belong to the same one of
$(-\infty,0)$ and $[0,\infty)$. Write their absolute values as
$0\le u<v$. Their separation gives $v-u\ge1$, whence
$v^2-u^2=(v-u)(v+u)\ge1$. Two numbers in $[m,m+1)$ have difference
strictly less than one, a contradiction. This includes equality in the
separation and bands meeting the closest point.
\end{proof}

\section{A sufficient parameter region}
Retain the source's coordinate maps, primary-factor dilation, adaptive
meshes, outer coloring, and digit product. Write
\[
 D=\lceil k^a\rceil,\quad M=\lceil k^m\rceil,\quad
 p=k^{-b},\quad h_0=D^2,\quad H=k^{-2},\quad
 \eta=(1000D)^{-1}.
\]
Let $q$ be the smallest power of the source prime $P\in(k,2k^2]$
at least $k^{ck/D}$. The heavy/light threshold remains $k/M$.
Here $m$ denotes an exponent, not a norm-band index or digit count.

\begin{proposition}\label{prop:region}
The parameter-dependent estimates in the released construction continue
to hold, using Lemma~\ref{lem:band}, if
\begin{gather*}
 a,b,c>0,\quad 3a<m<1,\quad b<a,\quad 3a<1-b,\\
 0<\delta<1,\quad 0<\gamma\le\delta/4,\quad
 \gamma<(1-\delta)/8,\quad 2c<\gamma b/2.
\end{gather*}
Consequently, conditional on the source's geometric and counting lemmas
with the parameter transfer proved below, for all sufficiently large $k$
and every integer $r\ge2$,
\[
 W_r(k)>k^{ck\lfloor\log_2r\rfloor}.
\]
The threshold may depend on the chosen parameters but not on $r$.
\end{proposition}
\begin{proof}
We track every scale that changes in the source Sections 2--7.
All parameters are fixed before $k$ grows.

The dilation estimate is
$\log\lambda\le D^2\log(2M)=O(k^{2a}\log k)$.
The inequalities $h_0<2M<k$ hold eventually since $2a<m<1$.
Thus $\lambda$ is coprime to $q$ as in the source, and each nonzero
cyclic progression has $k$ distinct terms. The prime-power rounding gives
\[
 ck\log k\le D\log q\le ck\log k+D\log(2k^2),
 \qquad \log q=O(k^{1-a}\log k).
\]
In particular the mesh proof still has
$R\le(D+1)\log q=O(k\log k)$ and $|V|=O(k^3\log k)$.
The width, neighborhood, and truncated-breakpoint bounds are unchanged.

The global-label count becomes
\[
 \log T_{\rm glob}=O(D(\log k+\log\lambda))
 =O(k^{3a}\log k)=o(M).
\]
This makes the first family of local-lemma weights
$kT_{\rm glob}\exp[-(\delta M-1)/32]=o(1)$.
The second family retains the bound
\[
 \sum_{h>D^2}(CDh)^{6D}e^{-h/64}=o(1),
\]
because $D\to\infty$ and
$6D\log(CDh)/h=O(\log D/D)=o(1)$ uniformly for $h>D^2$.
The source's absolute incident-weight and local-lemma criterion therefore
remain valid. The outer coloring still gives one quarter of each tested
set to each color.

For the drift and rational-path arguments it suffices that $M/k\to0$,
$MH\to0$, $MH/k\to0$, and $MH/(k\eta)\to0$.
Their exponents are respectively $m-1$, $m-2$, $m-3$, and $m+a-3$,
all negative in the stated region. The affinity condition $2kH<1$
is unchanged. In a long-period block the nonregular fraction is still
at most $4/1000+1/D<1/2$. Thus a block with a heavy index has at least
$h/8$ positions of each color. Such blocks cover at least
$(1-\delta)k-h$ positions, where $h\le2M=o(k)$.
They yield at least $((1-\delta)/8-o(1))k$ of each color.
The light case gives $\delta k/4$ of each color. Our conditions on
$\gamma$ imply the dichotomy with this new value of $\gamma$.
The source's short-period affine alternative is unchanged.

By Lemma~\ref{lem:band}, any rich progression requires at least
$\gamma k/2$ distinct flips to become one prescribed color. Hence
\[
 R_k\le2q^{2D}p^{\gamma k/2}
 \le\exp\bigl[(2c-\gamma b/2)k\log k
                  +2D\log(2k^2)+\log2\bigr]=o(1).
\]
Conflicting prescriptions have probability zero and cause no issue.
For affine progressions, the source's three-parameter arrangement count
gives
\[
 \log\max(1,T_{\rm aff})
 =O(k^{3a}\log k)+6\log q+O(\log k)
 =O((k^{3a}+k^{1-a})\log k)=o(k^{1-b}).
\]
Here the strict inequalities $3a<1-b$ and $b<a$ are essential.
Each consistent affine signature prescribes at least $k/2$ independent
bits, so
\[
 A_k\le2T_{\rm aff}(1-p)^{k/2}
 \le\exp[-pk/2+o(pk)]=o(1).
\]
Thus $R_k+A_k<1$ eventually. A cyclic coloring exists on
$N=q^D\ge k^{ck}$, and the source's unchanged digit-product argument
gives the claimed bound uniformly in $r$.
\end{proof}

\subsection{Explicit choices and scope of the improvement}
Three choices in Proposition~\ref{prop:region} are:
\begin{center}\small
\begin{tabular}{@{}rrrrrrr@{}}
\toprule
$a$&$m$&$b$&$\delta$&$\gamma$&$c$&Ratio to $10^{-5}$\\
\midrule
$1/10$&$1/2$&$1/20$&$1/10$&$1/40$&$1/4000$&$25$\\
$1/5$&$3/4$&$1/6$&$1/10$&$1/40$&$1/1000$&$100$\\
$1/4$&$4/5$&$247/1000$&$1/3$&$41/500$&$1/200$&$500$\\
\bottomrule
\end{tabular}
\end{center}
For the last row, the rich-event coefficient is exactly
$2c-\gamma b/2=-127/1000000$.
Both affine entropy power gaps equal $3/1000$; the light and long-block
color margins both equal $1/750$.
These small gaps can lead to a very large threshold. We claim no
practical finite coloring, finite van der Waerden record, or runtime gain.

In fact every fixed $c<1/192$ is admitted by this sufficient region:
take $\delta=1/3$, let $\gamma$ approach $1/12$ from below,
take $a=1/4$, let $b$ approach $1/4$ from below, and choose any
$m\in(3/4,1)$. Then $\gamma b/4$ approaches $1/192$.
This is a supremum within these retained estimates, not a universal
barrier on other improvements to the construction. The endpoint is not
asserted. Literature searches found no specific competing numerical
refinement, but this does not establish novelty or world-best status.

\section{An obstruction for the matrix profile inequalities}
The matrix paper defines a symmetric positive profile $P(a,b)$ with
$P(1,b)=b$, separate discrete concavity, and shifted tripling
$P(a,3h+a-1)\ge3P(a,h)$. Its rank envelope is
$P(a,b)\le(a+b-1)^{1/t}$. We exhibit an abstract profile compatible with
these conditions at $t=3/4$.

\begin{theorem}\label{thm:barrier}
For real $a,b\ge1$, set $u=\min(a,b)$, $v=\max(a,b)$ and
\[
 P_*(a,b)=\left(\frac{3u-1}{2}\right)^{1/3}
                      \left(v+\frac{u-1}{2}\right).
\]
Then $P_*$ is positive, symmetric, nondecreasing and separately concave,
$P_*(1,b)=b$, and it satisfies shifted tripling on positive integers.
Moreover
\[
 P_*(a,b)\le(a+b-1)^{4/3},\qquad
 P_*(a,a)=\left(\frac{3a-1}{2}\right)^{4/3}=\Theta(a^{4/3}).
\]
Thus these abstract constraints alone cannot force a larger diagonal
growth exponent, or exclude $t=3/4$.
\end{theorem}
\begin{proof}
Fix $a$. On $b\ge a$ the function is linear with positive slope.
On $1\le b\le a$, put $x=3b-1$, $A=3a-1$; it is
$(x/2)^{1/3}(2A+x)/6$. Its derivative in $b$ is positive and its
second derivative is
\[
 \frac{2^{2/3}(x-A)}{3x^{5/3}}\le0.
\]
The one-sided first derivatives match at $b=a$, both equal to
$((3a-1)/2)^{1/3}$. Concavity follows globally, and implies the required
discrete inequalities. Symmetry supplies the other coordinate.

If $h\ge a$, substitution gives
$P_*(a,3h+a-1)=3P_*(a,h)$ exactly. If $h<a$, put
$A=3a-1$, $B=3h-1$, and $z=(A/B)^{1/3}>1$.
After multiplication by a positive factor, shifted tripling becomes
\[
 z^4-2z^3+2z-1=(z-1)^3(z+1)\ge0.
\]
The boundary value and diagonal formula are direct substitutions.
Finally both $(3u-1)/2$ and $v+(u-1)/2$ are at most $a+b-1$,
proving the rank-envelope inequality.
\end{proof}
This is not a construction of a tensor character. It is an obstruction
to a proof using only the listed inequalities. New valid tensor
constraints, or a different construction, could still improve $9/4$.
The symbolic identity and exact rational cube comparisons were checked
by \texttt{audit\_math.py}; these supplement the displayed proof.

\section{Other tracks and verification limits}
A model-proposed search contracted edges and flipped diagonals in the
Sidorenko triple complex. It reproduced the released finite requirements
and tested 56,984 labeled candidates, finding no smaller object satisfying
the same discrete-refinement and three-cut certificate. This is neither
an exhaustive isomorphism-class search nor a minimality proof. A second
search tested 28,281 labeled candidates, including 2,518 discrete,
degree-eligible candidates, using rational-weighted cut distributions.
No feasible LP was returned. No impossibility certificate is claimed
from this floating-point feasibility screen. In particular these searches
do not exclude different pinning or singular-tail methods.

The unit-distance source was inspected for an explicit quantitative
target. Its proof uses asymptotic contradiction and successive limiting
arguments. No explicit improved exponent was derived in this pass.
The source manifest permits retrieval of the four inspected source trees.
Publication of this draft does not establish independent verification or
literature priority.
\end{document}
